% Regression: p3_eq10 — renders sum_{mu,i,c} [mu-bar fuses (a,b) -> c; i-bar fuses (c,x) -> y] [mirror].
% Formula: sum_{mu,i,c} [mu-bar fuses (a,b) -> c; i-bar fuses (c,x) -> y] [mirror]
% p3 — arXiv:2203.12563 eq (10): module associativity
%   sum_{mu,i,c} [mu-bar fuses (a,b) -> c; i-bar fuses (c,x) -> y] [mirror]
%     = sum_{i,j,z} [i-bar fuses (b,x) -> z; j-bar fuses (a,z) -> y] [mirror]
% The staircase geometry (one bar's combined leg entering another bar at
%   half-row height) has no grid form, so this is a free-tier sketch:
%   \tnput mpo boxes stand in for the fat fusion bars (no bar skin in the
%   free tier), and every wire is a black \tnjoin (no wire-hue key), so
%   the paper's red a, b, c wires and the black-vs-gray bar distinction
%   (module tensor mu vs action tensors i, j) are lost.
\documentclass[varwidth=19cm, border=6pt]{standalone}
\usepackage{amsmath, amssymb}
\usepackage{tenkz}
\begin{document}

\[
  \sum_{\mu,i,c}\;
  % term L1: mu fuses (a, b) -> c;  i fuses (c, x) -> y
  \begin{tenkzfree}
    \tnput[mpo]{m}{(14mm,4mm)}{\mu}
    \tnput[mpo]{i}{(26mm,-2mm)}{i}
    \tnjoin[label=a]{0mm,6.5mm}{m.north west}
    \tnjoin[label=b]{0mm,1.5mm}{m.south west}
    \tnjoin[label=c]{m.east}{i.north west}
    \tnjoin[label=x]{0mm,-4.5mm}{i.south west}
    \tnjoin[label=y]{i.east}{33mm,-2mm}
  \end{tenkzfree}
  \;
  % term L2 (mirror): i splits y -> (c, x);  mu splits c -> (a, b)
  \begin{tenkzfree}
    \tnput[mpo]{i}{(8mm,-2mm)}{i}
    \tnput[mpo]{m}{(20mm,4mm)}{\mu}
    \tnjoin[label=y]{0mm,-2mm}{i.west}
    \tnjoin[label=c]{i.north east}{m.south west}
    \tnjoin[label=x]{i.south east}{33mm,-4.5mm}
    \tnjoin[label=a]{m.north east}{33mm,6.5mm}
    \tnjoin[label=b]{m.east}{33mm,1.5mm}
  \end{tenkzfree}
  \;=\; \sum_{i,j,z}\;
  % term R1: i fuses (b, x) -> z;  j fuses (a, z) -> y
  \begin{tenkzfree}
    \tnput[mpo]{i}{(12mm,-2mm)}{i}
    \tnput[mpo]{j}{(24mm,3mm)}{j}
    \tnjoin[label=b]{0mm,0.5mm}{i.north west}
    \tnjoin[label=x]{0mm,-4.5mm}{i.south west}
    \tnjoin[label=a]{0mm,6.5mm}{j.north west}
    \tnjoin[label=z]{i.east}{j.south west}
    \tnjoin[label=y]{j.east}{31mm,3mm}
  \end{tenkzfree}
  \;
  % term R2 (mirror): j splits y -> (a, z);  i splits z -> (b, x)
  \begin{tenkzfree}
    \tnput[mpo]{j}{(8mm,3mm)}{j}
    \tnput[mpo]{i}{(20mm,-2mm)}{i}
    \tnjoin[label=y]{0mm,3mm}{j.west}
    \tnjoin[label=a]{j.north east}{31mm,6.5mm}
    \tnjoin[label=z]{j.south east}{i.north west}
    \tnjoin[label=b]{i.north east}{31mm,0.5mm}
    \tnjoin[label=x]{i.south east}{31mm,-4.5mm}
  \end{tenkzfree}
\]

\end{document}
